% =====================================================================
%  In-Depth Analysis of Homomorphic Encryption Libraries
%  Confidential Computing course project, Ben-Gurion University
%
%  Build:  cd report && tectonic final_report.tex && mv final_report.pdf ..
%  (or)    cd report && latexmk -pdf final_report.tex
%
%  Figures live in ../figures and are produced by src/plot_results.py.
% =====================================================================
\documentclass[10pt,a4paper]{article}

\usepackage[margin=2.0cm]{geometry}
\usepackage[T1]{fontenc}
\usepackage[utf8]{inputenc}
\usepackage{lmodern}
\usepackage{amsmath,amssymb}
\usepackage{booktabs}
\usepackage{graphicx}
\usepackage{float}
\usepackage[font=small,labelfont=bf,skip=4pt]{caption}
\usepackage[table]{xcolor}
\usepackage{microtype}
\usepackage{enumitem}
\usepackage[colorlinks=true,allcolors=blue!55!black]{hyperref}
\usepackage{titlesec}
\usepackage{abstract}

% Tighter, paper-like headings and lists.
\titlespacing*{\section}{0pt}{1.3ex plus .2ex}{0.7ex}
\titlespacing*{\subsection}{0pt}{1.0ex plus .2ex}{0.5ex}
\setlist{nosep,leftmargin=*}
\renewcommand{\abstractnamefont}{\normalfont\bfseries}
\renewcommand{\abstracttextfont}{\normalfont\small}
\setlength{\abovecaptionskip}{4pt}
\setlength{\belowcaptionskip}{2pt}

% Scale the wide, short chart images; at full width they waste vertical space.
\newcommand{\chart}[2]{%
  \begin{figure}[H]\centering
    \includegraphics[width=0.80\linewidth]{#1}
    \caption{#2}
  \end{figure}}

\newcommand{\ind}{IND-CPA\textsuperscript{D}}

\title{\vspace{-1.2cm}\textbf{In-Depth Analysis of Homomorphic Encryption Libraries}\\
       \large A measured comparison of TenSEAL and OpenFHE under CKKS}
\author{Itai Zloclzower \and Jacob Shemesh \and Tomer Ovadia \and Yoni Glickstein}
\date{Confidential Computing, Ben-Gurion University of the Negev \\ July 2026}

\begin{document}
\maketitle
\vspace{-0.6cm}

\begin{abstract}
\noindent
We compare two CKKS implementations --- TenSEAL (on Microsoft SEAL) and OpenFHE ---
against a plaintext baseline, across primitive operations, a vector-size sweep, and a
multi-stage computation over 1000 encrypted records. Two things distinguish this study
from a defaults-versus-defaults comparison. First, we re-run OpenFHE pinned to TenSEAL's
ring dimension and scaling factor, which shows that most of the apparent library gap is a
parameter gap: the multiplication penalty falls from $6.89\times$ to $3.16\times$, and the
geometric mean across six operations from $6.44\times$ to $3.05\times$ --- though decryption
remains $7.8\times$ and the matched arm is no longer 128-bit secure. Second, we measure
what the \ind{} decryption-noise defence costs, and find that only one of the two libraries
offers one at all: TenSEAL's decryption is bit-deterministic, while OpenFHE randomises by
default and supports calibrated flooding for $1.42\times$ on decryption. We also report a
cost the literature and our own earlier drafts routinely omit --- evaluation keys are
179.7\,MB against a 5.3\,MB encrypted cohort. This is a course project, and its findings
largely reproduce and illustrate known results; its contribution is measurement discipline,
including four claims from an earlier draft that did not survive re-examination and are
marked as retractions in place.
\end{abstract}

\vspace{0.2cm}

% =====================================================================
\section{Introduction and Goal}

Homomorphic Encryption (HE) allows computation directly on ciphertexts, so a cloud provider
can process data it cannot read. This project compares \textbf{TenSEAL} and \textbf{OpenFHE},
both implementing \textbf{CKKS}, against a plaintext baseline.

\textbf{Research question:} what is the practical overhead of computing under HE, and how
much of the measured difference between two libraries is a property of the \emph{libraries}
rather than of the \emph{parameters} we gave them?

The work has three parts: the cost of individual encrypted operations (\S\ref{sec:prim});
a multi-stage computation with realistic arithmetic structure --- a synthetic composite
score over 1000 encrypted records (\S\ref{sec:usecase}); and whether CKKS confidentiality
suffices for a scenario that releases decrypted results, which ours does (\S\ref{sec:sec}).

Four CKKS concepts recur. The \textbf{ring dimension} sets the security level and dominates
every operation's cost; the \textbf{scaling factor} sets precision; \textbf{multiplicative
depth} limits sequential multiplications; and \textbf{slot packing} lets one ciphertext hold
many values so that a single operation advances all of them.

This version follows an adversarial review of an earlier draft. Four of its claims did not
survive and are marked as retractions where they occur (\S\ref{sec:prim}, \S\ref{sec:mem},
\S\ref{sec:matched}, \S\ref{sec:results}); \S\ref{sec:limits} records what remains unfixed.

% =====================================================================
\section{Related Work}\label{sec:related}

Cross-library FHE benchmarking is an established genre, and we position this course project
inside it rather than claiming novelty.

\textbf{Library comparisons.} Sathya et al.~\cite{sathya2018} survey HE libraries and their
trade-offs. More recently, Faneela et al.~\cite{faneela2025} benchmark exactly the pair we
study --- SEAL and OpenFHE --- across BGV and CKKS on Linux and Windows, and conclude that
OpenFHE is the stronger choice across their settings. A 2025 benchmark study extends the
comparison to HElib and Lattigo~\cite{perfanalysis2025}. That our headline ordering differs
from Faneela et al.\ is itself informative and is the subject of \S\ref{sec:matched}: library
rankings in this area are highly sensitive to the parameters each library is given.

\textbf{Fair comparison is a known problem, not one we discovered.} Takeshita et
al.'s \emph{HEProfiler}~\cite{heprofiler2025} profiles four CKKS implementations and
explicitly ``devises methods for fair comparisons of these libraries, even with their widely
varied development strategies and library architectures'' --- the same difficulty our
matched-parameter arm addresses, treated far more thoroughly. Viand et al.'s
SoK~\cite{viand2021} surveys FHE compilers and the comparability problem across backends, and
the HEBench framework standardises workloads across libraries. Our \S\ref{sec:matched} is a
small, single-operating-point instance of a concern these works already articulate.

\textbf{Scheme and implementation.} CKKS is due to Cheon et al.~\cite{ckks2017}; both
libraries implement its full-RNS variant~\cite{rns2018}, where the modulus-chain and
rescaling behaviour we observe lives. The libraries are documented by Al Badawi et
al.~\cite{openfhe2022} and Benaissa et al.~\cite{tenseal2021}. OpenFHE's scaling techniques
--- whose $1.64\times$ cost and ${\sim}6\times$ precision effect we measure --- are introduced
by Kim et al.~\cite{kim2022}, which supplies the mechanism for a result we would otherwise
have reported as an unexplained tunable. Parameters follow the Homomorphic Encryption
Security Standard~\cite{hestandard}, the source of the 218/438-bit ceilings we use.

\textbf{Security of approximate FHE.} Li and Micciancio~\cite{limicciancio2021} defined
\ind{} and demonstrated key recovery against CKKS implementations including SEAL. Li et
al.~\cite{lmss2022} then obtained provable security by post-processing decryption with a
differential-privacy mechanism calibrated to \emph{worst-case} error growth; their result
\emph{reduces} the flooding required, since earlier approaches left CKKS roughly 8--16 bits of
precision. That worst-case qualifier is the hinge: Guo et al.~\cite{guo2024} break flooding
built on non-worst-case estimation, which is what OpenFHE's empirical estimator provides.
Bergamaschi et al.~\cite{pkc2025} analyse concrete security at reduced noise, and Alexandru
et al.~\cite{aahe2024} propose application-aware parameter selection.

\textbf{What we add, modestly.} Relative to the above: a matched-parameter ablation that
separates ring dimension from scaling technique at one operating point, and a measurement of
the \ind{} mitigation's cost placed alongside the performance numbers rather than discussed
separately. Neither is a new technique, and the security mechanism is documented by OpenFHE
itself. We contribute no framework and no cryptanalysis.

% =====================================================================
\section{Method}

\subsection{Configurations}

\begin{table}[H]\centering\small
\begin{tabular}{@{}lll@{}}
\toprule
 & \textbf{TenSEAL 0.3.15} & \textbf{OpenFHE 1.4.1} \\
\midrule
Configured by & modulus chain given directly & depth + scaling size; library picks the rest \\
Parameters & \texttt{poly\_mod=8192}, \texttt{[60,40,40,60]}, $2^{40}$ & \texttt{depth=3}, \texttt{scal\_mod=50}, \texttt{HEStd\_128} \\
Resulting ring dimension & 8192 & \textbf{16384} \\
\bottomrule
\end{tabular}
\end{table}

\noindent The two libraries are configured through different routes. That is convenient, and
it is also a serious comparability problem --- addressed in \S\ref{sec:matched}. All mandatory
operations were implemented for both, plus every optional extension: dot product, the score of
\S\ref{sec:usecase}, scaling tests, separately-timed key generation, the matched comparison,
and the \ind{} analysis.

\subsection{Measurement}\label{sec:method}

Timing an HE operation once is not reproducible. Every timing is a \textbf{mean over 1000
repetitions after 50 discarded warm-up calls}, except: key generation uses 50 repetitions, and
the matched-parameter and \ind{} experiments use 200. Warm-up matters because the first calls
pay one-off costs --- library loading, lazy NTT and key-cache construction, first-touch page
faults, frequency ramp-up --- that are not steady-state.

\textbf{Intervals are computed under dependence.} Consecutive timings share cache, allocator
and scheduler state, so they are not independent and $s/\sqrt{n}$ does not apply. On our
archived samples the lag-1 autocorrelation reaches \textbf{0.95} and the naive interval is
\textbf{2--6$\times$ too narrow}. Every interval here is a non-overlapping batch-means interval
(25 batches); the results files record both, plus the inflation factor.

\textbf{Stationarity uses three diagnostics.} A first-half/second-half comparison is a smoke
alarm, not a test: it is blind to a monotonic ramp (which halves its apparent size), to a
single spike, to disturbances symmetric about the midpoint, and to load present throughout. We
add a \textbf{Mann--Kendall} trend test and report lag-1 autocorrelation; a measurement is
flagged if either fires. A unit test documents the V-shaped case all three still miss. This
earned its place: our first run, on a shared node at load average 10.4 on 8 cores, showed
drifts to \textbf{73\%}, so everything here was re-measured on an \textbf{exclusive node},
single-threaded.

\textbf{Other precautions.} \texttt{tracemalloc} distorts short timings, so memory is a
separate pass; GC is disabled inside the timed loop as \texttt{timeit} does; each benchmark
asserts decrypted results against the plaintext reference before reporting timings --- a crash
detector for level exhaustion, deliberately loose, not a precision guarantee. 47 unit tests
cover the statistics, the drift blind spots, the score, and the padding invariant.

\textbf{Caveat on the plaintext denominators.} The baseline is pure Python, timed per call at
0.3--1.0\,\textmu s against ${\sim}130$\,ns of instrumentation overhead. Every ``$N\times$
slower'' figure is therefore \emph{relative to a pure-Python reference} and an upper bound on
the ratio against optimised plaintext. We added neither batched timing nor a NumPy baseline.

\textbf{Environment.} Exclusive SLURM node, Intel Xeon E5-2680 v3 @ 2.50\,GHz,
\texttt{OMP\_NUM\_THREADS=1}, Rocky Linux 9.7, Python 3.8.20. Versions, git commit and node are
recorded in every results file. Single-threading is deliberate --- thread count is not a
library property.

% =====================================================================
\section{Primitive Operations}\label{sec:prim}

Input \texttt{[1.0, 2.0, 3.0, 4.0, 5.0]}. ``infl.'' is how much wider each interval is than the
independence formula an earlier draft used.

\begin{table}[H]\centering\small
\begin{tabular}{@{}lrrrrrrrr@{}}
\toprule
\textbf{Operation} & \textbf{Plain (ms)} & \textbf{TenSEAL} & \textbf{CI $\pm$} & \textbf{infl.}
 & \textbf{OpenFHE} & \textbf{CI $\pm$} & \textbf{infl.} & \textbf{TS/OF vs plain} \\
\midrule
Encrypt      & ---      & 5.4914  & 0.045 & 4.6$\times$ & 15.1101 & 0.019 & 2.1$\times$ & --- \\
Add          & 0.000629 & 0.0961  & 0.000 & 1.4$\times$ & 0.3869  & 0.000 & 1.2$\times$ & 153$\times$ / 615$\times$ \\
Multiply     & 0.000550 & 4.7413  & 0.063 & 5.8$\times$ & 32.4827 & 0.033 & 1.6$\times$ & 8,616$\times$ / 59,025$\times$ \\
Sum          & 0.000320 & 10.5082 & 0.013 & 2.0$\times$ & 76.1546 & 0.083 & 2.2$\times$ & 32,832$\times$ / 237,943$\times$ \\
Average      & 0.000349 & 11.7961 & 0.157 & 6.0$\times$ & 81.3563 & 0.308 & 4.5$\times$ & 33,794$\times$ / 233,072$\times$ \\
Dot Product  & 0.000993 & 11.0764 & 0.014 & 2.3$\times$ & 96.6155 & 0.066 & 1.5$\times$ & 11,152$\times$ / 97,272$\times$ \\
Decrypt      & ---      & 1.4029  & 0.011 & 4.8$\times$ & 22.6702 & 0.053 & 2.9$\times$ & --- \\
\midrule
\textbf{Total, 5 comparable ops} & \textbf{0.002841} & \textbf{38.2180} & & & \textbf{286.9961} & & &
\textbf{13,450$\times$ / 101,003$\times$} \\
\bottomrule
\end{tabular}
\end{table}

Per-operation overhead varies by over two orders of magnitude --- $153\times$ for addition
against $32{,}832\times$ for summation --- so no single ``$N\times$ slower'' figure describes
HE. Addition is nearly free; rotation-heavy operations each need a $\log_2(\text{batch size})$
sequence of rotate-and-add steps, every one a key switch, costing 100--120$\times$ an addition.
TenSEAL leads on every operation, but \S\ref{sec:matched} shows most of that is parameters.

\paragraph{Two corrections to these totals.} The total now sums the \textbf{same five
operations} in all three columns; an earlier draft divided a seven-operation encrypted total
(including encrypt and decrypt, which plaintext lacks) by a five-operation plaintext one and
called it a slowdown, inflating the figures to $15{,}876\times$ and $114{,}299\times$. The
seven-operation session costs (45.11 / 324.78\,ms) are a deployment quantity, reported
separately.

More consequentially: an earlier draft argued that 1000 repetitions driving every interval
below a few percent made the measurements trustworthy. Our own samples disprove that twice. An
interval measures whether samples agree \emph{with each other} assuming independence --- the
contaminated run had narrow intervals \emph{and} 73\% drift. And the samples are not
independent even on a quiet machine: lag-1 autocorrelation reaches 0.95, so the naive interval
was too narrow by up to $6\times$. \textbf{Precision is not accuracy}, and this report made the
error it was warning about.

\subsection{Key Generation}
One-off setup, 50 repetitions: \textbf{TenSEAL 216.180\,ms} ($\pm$0.281), \textbf{OpenFHE
208.308\,ms} ($\pm$0.206). The only measurement where OpenFHE wins, and the two are within 4\%.
Setup dominates short sessions: TenSEAL's key generation costs nearly five times its entire
45\,ms benchmark of encrypted operations.

\subsection{Memory, and the Object We Forgot to Measure}\label{sec:mem}

An earlier draft presented \texttt{tracemalloc} figures --- all under 1\,KB --- as memory usage.
\textbf{That was misleading.} \texttt{tracemalloc} sees Python allocations only; both libraries
hold coefficients in native heaps, so those were wrapper-object sizes.

\begin{table}[H]\centering\small
\begin{tabular}{@{}lrr@{}}
\toprule
\textbf{Measure} & \textbf{TenSEAL} & \textbf{OpenFHE} \\
\midrule
\texttt{tracemalloc}, encrypt      & 0.59\,KB    & 0.42\,KB \\
\textbf{Resident per ciphertext}   & \textbf{359,465\,B} & \textbf{1,258,455\,B} \\
Serialized                          & 334,428\,B  & 1,312,467\,B \\
\bottomrule
\end{tabular}
\end{table}

That is \textbf{596$\times$} larger for TenSEAL and $2{,}947\times$ for OpenFHE --- two to four
orders of magnitude, not the six a previous draft claimed (a dropped kilobyte-to-byte
conversion, in the paragraph correcting a units error). Resident and serialized agree within
${\sim}7\%$. The two libraries are not size-comparable here: different ring dimension and tower
count, and SEAL compresses its serialisation where our OpenFHE path does not.

\paragraph{What actually crosses the network.} That paragraph led an earlier draft to conclude
ciphertext expansion ``is what crosses the network.'' Measuring the overlooked object reverses
it --- before the cloud can compute at all it needs the \textbf{evaluation keys}:

\begin{table}[H]\centering\small
\begin{tabular}{@{}lrr@{}}
\toprule
\textbf{At the deployed 1000-patient configuration} & \textbf{TenSEAL} & \textbf{OpenFHE} \\
\midrule
Encrypted cohort (5 features) & 5.27\,MB & 7.87\,MB \\
\textbf{Evaluation keys}      & \textbf{179.7\,MB} & \textbf{69.2\,MB} \\
Keys $\div$ cohort            & \textbf{34$\times$} & \textbf{8.8$\times$} \\
\bottomrule
\end{tabular}
\end{table}

They are one-off per context, which is why they are easy to omit --- and why omitting them
misleads any deployment that does not amortise a context over very many queries. Of everything
in this revision, this is the correction we would most want carried forward.

\subsection{Approximation Error}
Maximum absolute error, TenSEAL vs OpenFHE: Add $1.57\!\times\!10^{-9}$ vs
$4.09\!\times\!10^{-14}$; Multiply $3.35\!\times\!10^{-6}$ vs $1.96\!\times\!10^{-12}$; Sum
$4.60\!\times\!10^{-7}$ vs $7.11\!\times\!10^{-15}$; Average $3.10\!\times\!10^{-7}$ vs
$5.20\!\times\!10^{-14}$; Dot $5.19\!\times\!10^{-6}$ vs $3.27\!\times\!10^{-13}$. OpenFHE is
$10^{4}$--$10^{7}$ times more precise in this configuration; multiplication dominates the error
in both, as expected. All errors are negligible for aggregates meaningful to a few decimal
places. The \emph{cause} is not what we assumed --- see \S\ref{sec:matched}.

\subsection{Scaling with Vector Size}\label{sec:scaling}

End-to-end pipeline; context and keys built once per size and excluded.

\begin{table}[H]\centering\small
\begin{tabular}{@{}rrrrrrr@{}}
\toprule
\textbf{Size} & \textbf{Plain (ms)} & \textbf{TenSEAL} & \textbf{OpenFHE} & \textbf{TS vs pt} & \textbf{OF vs pt} & \textbf{OF/TS} \\
\midrule
5   & 0.001309 & 37.16  & 301.55 & 28,389$\times$ & 230,368$\times$ & 8.1$\times$ \\
10  & 0.001683 & 51.13  & 347.45 & 30,383$\times$ & 206,446$\times$ & 6.8$\times$ \\
50  & 0.004772 & 98.95  & 451.32 & 20,735$\times$ & 94,576$\times$  & 4.6$\times$ \\
100 & 0.008711 & 119.99 & 501.16 & 13,774$\times$ & 57,532$\times$  & 4.2$\times$ \\
500 & 0.044284 & 271.19 & 607.03 & 6,124$\times$  & 13,708$\times$  & 2.2$\times$ \\
\bottomrule
\end{tabular}
\end{table}

Over a $100\times$ data increase, plaintext grows $33.8\times$ but \textbf{TenSEAL only
$7.3\times$ and OpenFHE $2.0\times$}; per-element cost falls $13.7\times$ and $49.7\times$. The
mechanism is slot packing --- cost is set by ring dimension, not by how many slots hold data,
and OpenFHE held ring 16384 throughout. Only the rotation count, $\log_2(\text{batch size})$,
grows.

Two caveats. The plaintext baseline is \emph{itself} sublinear ($n^{0.77}$), because fixed
per-call interpreter overhead amortises the same way --- so this row is not independent evidence
for the CKKS mechanism. And the $n{=}5$ and $n{=}10$ plaintext points are both flagged by the
drift diagnostics ($-5.5\%$ and $-10.0\%$, relative SD above 50\%), so the largest overhead
figures rest on the least stable denominators in the study.

% =====================================================================
\subsection{Matched-Parameter Comparison}\label{sec:matched}

``TenSEAL is faster, OpenFHE is more precise'' is not supportable: OpenFHE was doing twice the
polynomial work per operation and keeping ten more bits per value. \texttt{SetRingDim} (with
\texttt{HEStd\_NotSet}) pins it to 8192 at $2^{40}$, matching TenSEAL. We also swept the
\emph{scaling technique}, which an earlier draft had silently confounded with the parameter
change.

\begin{table}[H]\centering\small
\begin{tabular}{@{}lrrrrrrr@{}}
\toprule
\textbf{Arm} & \textbf{Encr} & \textbf{Add} & \textbf{Mult} & \textbf{Sum} & \textbf{Dot} & \textbf{Decr} & \textbf{geo.} \\
\midrule
library-chosen (16384, $2^{50}$, FLEXIBLEAUTOEXT) & 2.52 & 4.02 & 6.89 & 7.24 & 8.70 & 16.28 & \textbf{6.44} \\
matched, \textbf{same technique} (FLEXIBLEAUTOEXT) & 1.20 & 2.12 & 3.16 & 3.25 & 3.97 & 7.83 & \textbf{3.05} \\
matched, FIXEDMANUAL                               & 1.02 & 1.77 & 1.93 & 2.41 & 3.09 & 6.80 & \textbf{2.37} \\
matched, FIXEDAUTO / FLEXIBLEAUTO                  & 1.02 & 1.75 & 1.93 & 2.40 & 3.09 & 6.82 & 2.37 \\
\bottomrule
\end{tabular}
\caption*{\small Ratio to TenSEAL, all six operations. An earlier draft showed only the two that converge best.}
\end{table}

\chart{../figures/chart_matched_comparison.png}{OpenFHE placed on TenSEAL's parameters. Lower is better in both panels.}

\begin{enumerate}[label=\textbf{\arabic*.}]
\item \textbf{Two variables, separated.} An earlier draft reported ``$6.89\times \to 1.93\times$
once parameters match'' and credited ring dimension alone --- while also changing the scaling
technique, which the same section showed is worth $1.64\times$. Holding technique fixed: ring
$16384\to8192$ gives \textbf{2.18$\times$}; FLEXIBLEAUTOEXT$\to$FIXEDMANUAL at ring 8192 gives
\textbf{1.64$\times$}; product $3.56\times$, exactly the observed closure. The honest
parameter-matched figure is \textbf{3.16$\times$}.

\item \textbf{The gap narrows everywhere but closes nowhere.} Geometric mean falls
$6.44\times \to 3.05\times$. Encryption converges to $1.20\times$; \textbf{decryption stays
$7.8\times$} --- the largest residual, on the operation \S\ref{sec:sec} makes security-critical,
and one the earlier table omitted. The summary is a range: \textbf{1.2--7.8$\times$ by
operation}.

\item \textbf{Precision is not explained by scale alone.} At identical $2^{40}$ OpenFHE stays
275--1{,}890$\times$ more precise. Calling that an \emph{implementation} difference is too
strong: the arms also differ in key-switching auxiliary budget and batch size, neither varied,
and technique alone moves precision ${\sim}6\times$. The mechanism is not established.

\item \textbf{The matched arm is not 128-bit secure --- a retraction.} An earlier draft claimed
it was, comparing OpenFHE's reported modulus (181--200 bits) against SEAL's 218-bit ceiling.
Different quantities: SEAL's bound includes its special prime, OpenFHE's figure omits its
hybrid auxiliaries. The decisive test is to ask the library --- at \texttt{HEStd\_128\_classic}
with these parameters \textbf{OpenFHE selects ring 16384 and refuses 8192}, for all four
techniques. The experiment now records that refusal. Parameter parity cost security parity.
\end{enumerate}

Still a \textbf{trade-off curve, not a ranking} --- but weaker than first claimed. This also
offers a reading of why Faneela et al.~\cite{faneela2025} rank OpenFHE above SEAL where we rank
TenSEAL above OpenFHE: at defaults the two libraries are not running comparable parameters, and
the ordering is a property of the configuration route as much as of the code.

% =====================================================================
\section{Applied Use Case: An Encrypted Composite Score}\label{sec:usecase}

\noindent\fbox{\parbox{0.97\linewidth}{\small \textbf{Disclaimer.} The score below is a
\textbf{synthetic construct for demonstration only}. It is \textbf{not} a validated clinical
model, is not derived from medical literature, and must not be read as an assessment of patient
risk. The cohort is synthetic. Its purpose is to give the benchmark a computation with realistic
\emph{structure} --- several inputs, mixed arithmetic, genuine multiplicative depth.}}

\subsection{Threat Model, and Why This Workload Is a Proxy}\label{sec:threat}

Earlier drafts described the scenario without naming an adversary. \textbf{Assets:} individual
feature values; the released cohort mean; arguably the formula. \textbf{Adversary:} a
semi-honest cloud that follows the protocol but inspects everything --- ciphertexts, circuit,
timings, released outputs. \textbf{Trust:} the hospital holds plaintext and the secret key.

\textbf{What this does not protect.} The circuit is public to the cloud (weights, structure and
ranges are plaintext constants it evaluates). Cohort size and query pattern leak. The
\emph{released aggregate} is a disclosure channel HE does not touch --- output privacy, via
differential privacy over the decrypted value, is a different problem from \ind{} and we address
neither. And there is \textbf{no integrity}: a malicious or buggy cloud can return a ciphertext
decrypting to anything, undetected. Our accuracy tables check against a plaintext reference the
hospital computed itself, which a real deployment could not do.

\textbf{Why this workload is a proxy.} We compare HE against plaintext only, which answers
``what does HE cost'', not ``should you use HE''. The real alternatives --- TEEs, secure
multi-party computation, differential privacy alone, or not outsourcing --- are outside our
scope and we benchmark none. And \S\ref{sec:results} measures the plaintext score for all 1000
records at 0.075\,ms on 40\,KB of data: \textbf{no one would outsource that}, and the client
spends more CPU encrypting and decrypting than computing the score itself. The scenario earns
its place because its \emph{arithmetic structure} matches workloads where outsourcing is
justified.

\subsection{The Score}
Each feature is min-max normalised against a \emph{public} reference range --- age 18--90,
systolic BP 90--180, BMI 18--40, cholesterol 120--280, glucose 70--200:
\begin{align*}
\text{score} = 100 \cdot \big(&0.15\,a + 0.20\,b + 0.15\,m + 0.15\,c + 0.15\,g \\
                              &+ 0.10\,b\,g + 0.05\,m\,c + 0.05\,b^2\big)
\end{align*}
where each symbol is the normalised feature. The eight weights sum to 1.0. The first five terms
need only multiplication by \emph{public constants}. The last three make this a real HE
benchmark: two \textbf{interaction terms} and one \textbf{quadratic term}, each requiring
\textbf{ciphertext $\times$ ciphertext} multiplication, which consumes depth and forces
relinearization. The weights were chosen so the terms sum to 1.0 and so that the last three
force ciphertext-by-ciphertext products; they are not fitted to any outcome.

Ranges are public deliberately: deriving min/max from the data would need ciphertext
comparison, which CKKS does not support natively. \textbf{Cohort}: 1000 records, fixed seed;
age uniform, the rest normal-then-clipped, giving scores of mean 35.63, SD 8.45, range
9.36--68.90.

\subsection{Encrypted Implementation}
\textbf{Packing}: one feature per ciphertext, one slot per record --- five ciphertexts of 1000
values, so one multiplication advances all 1000. \textbf{Padding}: vectors are padded to 1024
slots (OpenFHE requires a power of two) with \emph{each feature's range minimum}, so they
normalise to 0 and drop out of every term and the cohort sum. Padding with zeros instead gives
those slots a nonzero normalised value of $-\min/(\max-\min)$ and silently corrupts the mean; a
unit test pins this. \textbf{Depth}: normalisation (1) + interaction/quadratic (1) + term
weights (1) + cohort mean (1) $=$ \textbf{4 levels}.

\subsection{Depth, Not Data Volume, Drives Cost}
The \S\ref{sec:prim} parameters \textbf{cannot run this circuit}. TenSEAL at
\texttt{poly=8192} fails with \texttt{scale out of bounds}, and any longer chain is rejected
outright; \texttt{poly=16384} with 4 levels works. OpenFHE at depth 3 evaluates the score but
fails on the cohort mean; depth 4 works at ring 16384; depth 5 works but pushes the ring to
32768. The rejection is not a quirk: SEAL's published table caps the coefficient modulus at
\textbf{218 bits for ring 8192} and 438 for 16384, and $[60,40,40,60]=200$ fits while
$[60,40,40,40,60]=240$ does not.

So a modest increase in circuit complexity forces new cryptographic parameters, and asking for
more depth than needed is an easy and costly mistake --- our first implementation requested
depth 5, doubling the ring for no benefit. Two caveats we previously omitted: this holds at a
\emph{fixed scaling factor} (the constraint is depth $\times$ scale), and depth 4 is minimal for
\emph{our formulation}, not the circuit --- folding term weights into the normalisation
constants and doing $\times 1/N$ after decryption would give depth 2.

\subsection{Results}\label{sec:results}

\begin{table}[H]\centering\small
\begin{tabular}{@{}llrrr@{}}
\toprule
\textbf{Stage} & \textbf{Runs at} & \textbf{Plaintext (ms)} & \textbf{TenSEAL (ms)} & \textbf{OpenFHE (ms)} \\
\midrule
Encrypt 5 features & Hospital & ---     & 76.784  & 84.650  \\
Score evaluation   & Cloud    & 0.0593  & 89.249  & 163.181 \\
Cohort mean        & Cloud    & 0.0155  & 46.104  & 146.054 \\
Decrypt scores     & Hospital & ---     & 2.186   & 17.457  \\
\midrule
\textbf{Pipeline total}       & & \textbf{0.0748} & \textbf{214.323} & \textbf{411.341} \\
\textbf{Overhead vs plaintext}& & \textbf{1$\times$} & \textbf{2,865$\times$} & \textbf{5,499$\times$} \\
\textbf{Per record}           & & \textbf{0.000075} & \textbf{0.2143} & \textbf{0.4113} \\
Cloud-side only               & Cloud & 0.0748 & 135.353 & 309.235 \\
\bottomrule
\end{tabular}
\end{table}

Accuracy: maximum absolute error $4.881\!\times\!10^{-4}$ (TenSEAL) and
$1.976\!\times\!10^{-10}$ (OpenFHE) score points; encrypted cohort mean 35.627391 and 35.626948
against a plaintext 35.626948.

\chart{../figures/chart_risk_score.png}{Encrypted score over 1000 records: stage cost, cohort distribution, accuracy.}

\paragraph{The overhead figures are a correction.} An earlier draft reported $3{,}614\times$ and
$6{,}936\times$, using a plaintext denominator of 0.0593\,ms --- the score stage alone ---
against an encrypted numerator summing four stages. Worse, the two plaintext stages
\emph{overlapped}: the cohort-mean stage recomputed the scores, so it contained the score stage
and the two could not be summed either. Stages are now disjoint and the ratios fall ${\sim}21\%$.

The scenario runs in \textbf{214\,ms} (TenSEAL) or 411\,ms (OpenFHE) for 1000 records. Overhead
is well below the isolated primitives ($2{,}865\times$ against $13{,}450\times$): nothing got
faster, the ciphertexts got fuller --- though still only one-eighth full, so the amortisation
argument is made below saturation.

\paragraph{What the per-record number omits.} It says computation is not the bottleneck; it does
not say the system is deployable. Excluded: key generation; the \textbf{evaluation keys} at
179.7\,MB and 69.2\,MB (\S\ref{sec:mem}); transport of 5.3\,MB; key custody and rotation. Client
and server also share one process holding the secret key, so ``the cloud never sees plaintext''
is a property of the simulation.

% =====================================================================
\section{Security: Is Encrypting Enough?}\label{sec:sec}

CKKS is IND-CPA secure. IND-CPA says nothing about releasing \emph{decryption results} --- and
our scenario releases them. A CKKS decryption returns the message \emph{plus an error term that
depends on the secret key}; an adversary that obtains approximate decryptions of ciphertexts
whose exact plaintexts it can predict can solve for the key~\cite{limicciancio2021}. We
deliberately do not quote a query count.

\subsection{What the Libraries Do, and What the Defence Costs}

\S\ref{sec:sec} runs on a \textbf{64-record} cohort (batch 64, depth 4, ring 16384, 200
repetitions), not \S\ref{sec:usecase}'s 1000; the noise estimate is circuit- and data-dependent
and does not transfer unchanged. These arms are \textbf{not matched to each other} --- the
criticism \S\ref{sec:matched} levels at the performance comparison applies here too.

The observable signature of a decryption-noise defence is \textbf{randomization}: plain CKKS
decryption is deterministic, so one fixed ciphertext must decrypt identically every time. Over
32 trials: \textbf{TenSEAL} --- not randomized, spread \textbf{exactly 0}; \textbf{OpenFHE
\texttt{FIXED\_NOISE\_DECRYPT}} (default) --- randomized, $5.71\!\times\!10^{-11}$;
\textbf{OpenFHE \texttt{NOISE\_FLOODING\_DECRYPT}} --- randomized, $1.28\!\times\!10^{-13}$.

\textbf{In the version tested TenSEAL's decryption is bit-deterministic}, so the released value
carries exactly the key-dependent error the attack exploits, and no setting changes it. OpenFHE
randomizes by default --- which surprised us.

\textbf{The table contains an inversion we cannot explain away.} The ``calibrated'' arm adds
${\sim}450\times$ \emph{less} noise than the default it supposedly improves on, and is
correspondingly more accurate. The consistent reading is that the default applies a conservative
constant while flooding calibrates to this circuit's measured error, which for so shallow a
circuit is smaller. If so, the two-pass arm is not straightforwardly stronger. We did not
establish which mode is more secure.

Flooding costs \textbf{1.42$\times$ on decryption and nothing elsewhere} (encrypt, multiply and
sum all $1.00\times$), plus a full extra pass for estimation.

\chart{../figures/chart_ind_cpad.png}{Cost of the flooding defence, and which configurations randomize the released value.}

\subsection{Why We Do Not Claim This Makes It Secure}

OpenFHE estimates circuit noise \textbf{empirically}, from precision loss in the imaginary slots
of a trial decryption --- an \textbf{average-case} estimate. Guo et al.~\cite{guo2024} show that
flooding built on non-worst-case estimation remains vulnerable to key recovery, explicitly
including deployments implementing the CRYPTO 2022 bounds. Average-case analysis assumes
independent input ciphertexts; our circuit combines correlated ones.

\textbf{So \S\ref{sec:sec} measures the cost of a partial defence, not the cost of security.}
Randomized decryption shows a defence is \emph{active}, not \emph{sufficient}; we neither
mounted the attack nor verified a bound. The same restraint applies in the other direction, and
an earlier draft did not apply it: determinism establishes that TenSEAL offers \textbf{no
mitigation}, not that our scenario is exploitable.

% =====================================================================
\section{Usability}

\begin{table}[H]\centering\small
\begin{tabular}{@{}lll@{}}
\toprule
\textbf{Aspect} & \textbf{TenSEAL} & \textbf{OpenFHE} \\
\midrule
Installation & \texttt{pip install}, all platforms & \textbf{Linux only}, one CPython version \\
API & High-level (\texttt{enc + enc}) & Low-level (\texttt{cc.EvalAdd(ct, ct)}) \\
Parameter control & Modulus chain explicit & Depth given; ring chosen by library \\
Scaling technique & Not exposed & Four options, real trade-off (\S\ref{sec:matched}) \\
Decryption-noise defence & \textbf{None} & Default fixed noise; opt-in flooding \\
\bottomrule
\end{tabular}
\end{table}

TenSEAL is far easier for prototyping. OpenFHE exposes controls that matter a great deal
(\S\ref{sec:matched}, \S\ref{sec:sec}) and that TenSEAL does not offer at all. One genuine
advantage of its style: declaring depth up front makes exceeding it fail immediately, where
TenSEAL's failure surfaces later as \texttt{scale out of bounds}.

% =====================================================================
\section{Limitations}\label{sec:limits}

Several were identified by adversarial review of an earlier draft and are recorded rather than
fixed, because fixing them needs measurement we did not do.

\textbf{Measurement.} (1) \textbf{Arms are not interleaved or replicated} --- one process, fixed
order, so every interval is \emph{within-run} and run-to-run variance (often dominant;
\cite{mytkowicz2009}) is unmeasured; two runs of the same OpenFHE configuration in our own
artifacts differ 10\% on encryption. (2) \textbf{Plaintext denominators are interpreter
measurements} with relative SD above 50\%; ratios against them are upper bounds relative to
optimised plaintext. (3) \textbf{Warm-up is short for TenSEAL}, whose transient runs a few
hundred calls, biasing its means slightly high --- an asymmetric effect. (4) \textbf{The
stationarity checks miss disturbances symmetric about the midpoint.} (5) \textbf{One node, one
CPU generation, single-threaded} --- AVX2 but not AVX-512, and the two libraries have different
vectorisation backends, so even the ratios may shift on newer hardware.

\textbf{Cryptographic scope.} (6) \textbf{The matched arm is not 128-bit secure} (see
\S\ref{sec:matched}). (7) \textbf{The two ``multiply'' operations differ}: TenSEAL relinearises
and rescales, OpenFHE defers the rescale, so $3.16\times$ is a lower bound. (8) \textbf{Depth 4
is minimal for our formulation, not the circuit.} (9) \textbf{``Depth drives ring dimension''
holds at fixed scale.} (10) \textbf{Ciphertext-size comparison is parameter-determined}, and
SEAL compresses where our OpenFHE path does not.

\textbf{Scenario.} (11) Semi-honest adversary only; \textbf{no integrity}. (12) \textbf{No
alternatives benchmarked} --- no TEE, MPC or DP baseline. (13) \textbf{The per-record figure
excludes most of a deployment}, and client and server share one process. (14) \textbf{Output
privacy and regulation are untreated.} (15) \textbf{Determinism is an observational bound} at
$n{=}32$ trials. (16) \textbf{The score is synthetic throughout}, with independently drawn
features where real measures correlate. (17) \textbf{No bootstrapping.}

% =====================================================================
\section{Conclusions}

\begin{enumerate}[label=\textbf{\arabic*.}]
\item \textbf{Cost depends on the operation and on how full the ciphertexts are.} Overhead ran
from $153\times$ (addition) to $32{,}832\times$ (summation) on a 5-element vector, falling to
$2{,}865\times$ for the whole pipeline once 1000 values shared each ciphertext --- still only
one-eighth full. A slowdown quoted without a fill fraction is close to uninterpretable.

\item \textbf{Most of the library gap was a parameter gap --- but less than we first said.}
Matched, and holding the scaling technique fixed, OpenFHE's multiply penalty falls
$6.89\times \to 3.16\times$, not the $1.93\times$ an earlier draft got by changing two variables
at once. Decryption still trails $7.8\times$. And the matched arm is \textbf{not} 128-bit secure.

\item \textbf{Depth, not data volume, drove cost in the regime measured} --- at fixed scaling
factor, and only while data fits the available slots.

\item \textbf{Rigorous measurement changed the numbers, then changed them again.} Contaminated
runs showed 73\% drift behind narrow intervals; the intervals were \emph{also} up to $6\times$
too narrow on a quiet machine, because consecutive timings are autocorrelated. Both are the same
error --- treating an interval as evidence of validity --- and this report made it twice.

\item \textbf{The largest object crossing the network is the one we did not measure.} Evaluation
keys are 179.7\,MB against a 5.3\,MB encrypted cohort.

\item \textbf{Encryption is not the whole security question.} Our scenario releases decrypted
aggregates, exactly where CKKS is not \ind{} secure. TenSEAL offers no mitigation; OpenFHE
randomises by default and offers calibrated flooding for $1.42\times$ on decryption --- but this
is a partial defence at best, and we could not explain why the calibrated arm adds less noise
than the default.

\item \textbf{Who should act.} Engineers choosing a CKKS library for a pipeline that
\emph{releases decrypted values} should weigh \S\ref{sec:sec} above \S\ref{sec:prim}. Anyone
quoting HE benchmarks should state fill fraction and parameters, which move headlines by an
order of magnitude --- and, as \S\ref{sec:matched} suggests, can reverse a library ranking.
\end{enumerate}

% =====================================================================
\begin{thebibliography}{99}\small
\bibitem{ckks2017} J.\,H. Cheon, A. Kim, M. Kim, Y. Song. Homomorphic encryption for arithmetic of approximate numbers. \emph{ASIACRYPT}, 2017.
\bibitem{rns2018} J.\,H. Cheon, K. Han, A. Kim, M. Kim, Y. Song. A full RNS variant of approximate homomorphic encryption. \emph{SAC}, 2018.
\bibitem{openfhe2022} A. Al Badawi et al. OpenFHE: open-source fully homomorphic encryption library. \emph{WAHC}, 2022.
\bibitem{tenseal2021} A. Benaissa, B. Retiat, B. Cebere, A.\,E. Belfedhal. TenSEAL: a library for encrypted tensor operations. 2021.
\bibitem{kim2022} A. Kim, A. Papadimitriou, Y. Polyakov. Approximate homomorphic encryption with reduced approximation error. \emph{CT-RSA}, 2022.
\bibitem{hestandard} M. Albrecht et al. Homomorphic Encryption Security Standard. HomomorphicEncryption.org, 2018.
\bibitem{viand2021} A. Viand, P. Jattke, A. Hithnawi. SoK: fully homomorphic encryption compilers. \emph{IEEE S\&P}, 2021.
\bibitem{heprofiler2025} J. Takeshita, N. Koirala, C. McKechney, T. Jung. HEProfiler: an in-depth profiler of approximate homomorphic encryption libraries. \emph{J. Cryptographic Engineering}, 2025.
\bibitem{faneela2025} Faneela, J. Ahmad, B. Ghaleb, S.\,U. Jan, W.\,J. Buchanan. Cross-platform benchmarking of the FHE libraries: novel insights into SEAL and OpenFHE. arXiv:2503.11216, 2025.
\bibitem{perfanalysis2025} Performance analysis of leading homomorphic encryption libraries: a benchmark study of SEAL, HElib, OpenFHE and Lattigo. \emph{ICCSAIDE}, 2025.
\bibitem{sathya2018} S.\,S. Sathya, P. Vepakomma, R. Raskar, R. Ramachandra, S. Bhattacharya. A review of homomorphic encryption libraries for secure computation. arXiv:1812.02428, 2018.
\bibitem{limicciancio2021} B. Li, D. Micciancio. On the security of homomorphic encryption on approximate numbers. \emph{EUROCRYPT}, 2021.
\bibitem{lmss2022} B. Li, D. Micciancio, M. Schultz, J. Sorrell. Securing approximate homomorphic encryption using differential privacy. \emph{CRYPTO}, 2022.
\bibitem{guo2024} Q. Guo, D. Nabokov, E. Suvanto, T. Johansson. Key recovery attacks on approximate homomorphic encryption with non-worst-case noise flooding countermeasures. \emph{USENIX Security}, 2024.
\bibitem{pkc2025} F. Bergamaschi et al. Revisiting the security of approximate FHE with noise-flooding countermeasures. \emph{PKC}, 2025.
\bibitem{aahe2024} A. Alexandru, A. Al Badawi, D. Micciancio, Y. Polyakov. Application-aware approximate homomorphic encryption. ePrint 2024/203.
\bibitem{mytkowicz2009} T. Mytkowicz, A. Diwan, M. Hauswirth, P.\,F. Sweeney. Producing wrong data without doing anything obviously wrong! \emph{ASPLOS}, 2009.
\end{thebibliography}

\vspace{0.2cm}
\noindent\rule{\linewidth}{0.4pt}\\
{\footnotesize Code, notebook and charts accompany this report; every figure is regenerable with
\texttt{./run\_all\_benchmarks.sh}. Software measured: TenSEAL 0.3.15, OpenFHE 1.4.1.}

\end{document}
